% Cote 119: offprints, reports, courses, CV and letters (typed and printed part): archive cote 119, site page 176 (PDF page 177).
% typescript. Draft compared: none.
% AI-transcribed by Claude Opus 5.5 (claude-opus-5-5, adaptive thinking, effort high), Message Batches API, batch r119_full (msgbatch_01DwZNm5Y7YsKBWyAN4sFs8w), 2026-11-13; not yet checked by eye.
\scanpage{119/176}
\mspage{- 15 -}

que le développement en série entière autour d'un point $Z_0$ du domaine $U$
de $f$ était valable dans le plus grand disque ouvert de centre $Z_0$ contenu
dans $U$ ce qui n'était pas évident à partir de la définition des fonctions
analytiques.

(5.7.) \uline{Théorème de Morera} (réciproque de 4.3.)

Soit $f$ une fonction continue dans \unsure{uncouvert|un ouvert} $U$ du plan. Si la forme
différentielle $f(Z)\,dZ$ est fermée alors $f$ est holomorphe dans $D$.

En effet $f$ admet localement une primitive $g$. Cette primitive est
holomorphe (elle a une dérivée $f$) et d'après (5.5.) $g$ est analytique donc
$f$ est aussi analytique C.Q.F.D.

En particulier si $f$ est continue dans $U$ et holomorphe en tout point de $U$
sauf éventuellement en des points isolés $f$ est holomorphe dans $U$ tout
entier card'après (4.4.) $f(Z)\,dZ$ est fermée et (5.7.) nous donne la
conclusion.

6 . \uline{Homotopie de chemins}

Une conséquence des propositions 3.4. et 3.6. est que les disques
ouverts du plan possèdent la propriété suivante :

6.1. Toute forme différentielle fermée sur un disque ouvert admet une
primitive sur ce disque ouvert.

On va voir que l\unsure{e}s disque sont loin d'être les seuls ouverts à posséder cette
propriété. En fait les ouverts possédant cette propriété sont les ouverts
``simplement connexe'' c'est-à-dire ceux qui n'ont pas de ``trous''. Nous allons
préciser cette notion vague en donnant un sens clair à la notion de chemin qui
se déforme continûment.

Pour simplifier on ne va considérer que des chemins paramétrés par le segment
$I = \hand{[}0,1\hand{]}$ . (on pourra de toutes façons se ramener à ce cas par un changement
de paramètre).

\margin{176}
